% chapters/14/checklist.tex
\bagian{Checklist Pencapaian Kompetensi}

Bab ini menuntaskan \SUBCPMK{4.1}, yaitu bab pertama pada \CPMK{4}. Centang
setiap butir berikut hanya apabila Anda benar-benar telah melakukannya atau
dapat menjelaskannya.

\begin{checklistbox}
\begin{checklist}
  \item Saya dapat menjelaskan tiga langkah jalan pengolahan aset pada Unity.
  \item {\bfseries (Indikator a)} Saya dapat mengatur \istilah{Sprite Mode},
        \istilah{Pixels Per Unit}, \istilah{Filter Mode}, dan
        \istilah{Compression} pada aset saya, dan menjelaskan alasan setiap
        pilihannya.
  \item Seluruh aset saya memakai nilai \istilah{Pixels Per Unit} yang sama.
  \item Saya dapat menghitung tinggi sebuah gambar pada scene dari ukuran
        berkasnya dan nilai \istilah{Pixels Per Unit}.
  \item Saya dapat menjelaskan mengapa ukuran berkas gambar tidak boleh diubah
        di luar Unity.
  \item {\bfseries (Indikator b)} Saya telah menyusun \istilah{Animator} dengan
        keadaan diam dan berlari beserta transisinya.
  \item Saya dapat menjelaskan perbedaan angka pecahan, nilai benar-salah, dan
        pemicu pada parameter \istilah{Animator}.
  \item Animasi pada scene saya berganti menurut gerak pemain.
  \item Saya dapat menjelaskan mengapa kesalahan penamaan parameter tidak
        menimbulkan pesan kesalahan, dan bagaimana mencegahnya.
  \item {\bfseries (Indikator c)} Saya telah menyusun \istilah{Canvas} yang
        menampilkan skor dan nyawa.
  \item Antarmuka saya benar-benar berubah ketika skor bertambah atau nyawa
        berkurang.
  \item Saya dapat menjelaskan mengapa keadaan permainan disimpan pada kode,
        bukan pada tampilannya.
  \item Saya telah membandingkan scene sebelum dan sesudah penggantian aset,
        dan memeriksa bentuk batas bendanya.
  \item Daftar aset saya memuat sumber dan lisensi setiap aset yang dipakai.
\end{checklist}
\end{checklistbox}
